\chapter{Deployment, restart and recovery}
\label{ch:deploy}

\section{Configuration management is the bulk of the work}

Not scheduling. Every machine needs an install, a configuration, a token, and
then repeated revisions of the owner policy as the policy is refined.

\begin{table}[H]
\centering
\begin{tabularx}{\textwidth}{L{0.20\textwidth} Y}
\toprule
\textbf{Item} & \textbf{Decision} \\
\midrule
Channel & The existing authenticated ssh fan-out to every host \\
Layout & \file{/etc/condor/config.d/} drop-ins, rendered from the repository \\
Loop & render $\rightarrow$ push $\rightarrow$ \kn{condor\_reconfig} \\
Slot changes & Require \kn{condor\_restart} of the startd, not \kn{condor\_reconfig} \\
Verification & \kn{condor\_config\_val -v <KNOB>} --- never trust the file, \file{config.d} precedence bites \\
\bottomrule
\end{tabularx}
\caption{Deployment mechanics.}
\end{table}

\section{Adding a worker}

\begin{enumerate}
  \item Install the package.
  \item Drop in the worker configuration: \kn{CONDOR\_HOST},
        \kn{DEFAULT\_DOMAIN\_NAME}, \kn{UID\_DOMAIN}, donation limits, the owner
        policy, and the \kn{STARTD\_CRON} stanza.
  \item Install \kn{condor-owner-session} to \file{/usr/local/bin} (mode 755).
  \item Install a \textbf{daemon} token to \file{/etc/condor/tokens.d/} (600, root).
  \item Restart, then confirm the machine appears in \kn{condor\_status}
        \textbf{and} that a job actually lands on it.
\end{enumerate}

Step 4 is the one that is easy to miss, for the reason given in
Chapter~\ref{ch:security}. Issue it on the entry node with the \emph{daemon}
identity:

\begin{lstlisting}[style=term]
sudo condor_token_create -identity condor@<pool>.internal
\end{lstlisting}

Move it straight from entry node to worker and destroy every intermediate copy.
It is a bearer credential with no machine binding: anything that reads the file
can act as a pool daemon.

Step 5's second clause is not redundant with its first: appearing in
\kn{condor\_status} and accepting work are independent properties, and
Chapter~\ref{ch:failure} catalogues several ways to have the first without the
second.

\section{Adding a user PC}

A user PC is a worker that also submits. Same procedure, from the same template,
plus three things:

\begin{itemize}
  \item \textbf{Set \kn{RANK}} to that person's identity. This is the whole
        adoption argument, and it is the only thing they get for joining that a
        donated-only worker does not.
  \item \textbf{Pin \kn{NETWORK\_INTERFACE}} if the machine carries a VPN or
        any interface the entry node cannot reach (\S\ref{sec:vpn}).
  \item \textbf{No schedd.} Submission is remote; the local install exists to
        donate capacity.
\end{itemize}

The pool and naming block is identical to a plain worker's, so this file
\emph{replaces} any earlier submit-client configuration rather than sitting
beside it. Leaving both in place means two files defining
\kn{CONDOR\_HOST}, resolved by \file{config.d} precedence rather than by
intent.

\section{What survives a reboot}

Verified by rebooting the entry node.

Everything survives, provided the services are enabled: configuration drop-ins,
the \file{/usr/local/bin} helpers, the pool signing key, and --- importantly ---
the \textbf{job queue}. Completed jobs were still in \file{job\_queue.log}
afterwards.

The pool also degrades gracefully rather than failing. With the worker absent it
ran a six-job cluster on the entry node's two donated cores and queued the
remainder.

\section{A central-manager restart costs about two update intervals}
\label{sec:cmrestart}

Workers do \textbf{not} come back promptly. The observed sequence:

\begin{lstlisting}[style=term]
17:19:41  collector starts as a NEW process
17:19:45  worker MASTER update  -> "SECMAN: Server rejected our session id"
17:24:00  worker STARTD update  -> "SECMAN: Server rejected our session id"
17:29:00  worker STARTD update  -> succeeds, re-registers
\end{lstlisting}

Each daemon holds a cached security session tied to the \emph{old} collector
process. It spends one full update cycle discovering that the session is dead,
invalidates it, and only re-authenticates on the \textbf{next} cycle --- about
ten minutes at the default 300-second \kn{UPDATE\_INTERVAL}.

\begin{trap}
It self-heals with no intervention, but \textbf{nothing in \kn{condor\_status}
explains it} --- the machine is simply missing. With many workers they all
disappear together after any entry-node restart, which looks like a catastrophe
and is not.
\end{trap}

The evidence lives in \file{CollectorLog} (\kn{DC\_AUTHENTICATE: attempt to open
invalid session}) and in the worker's own log (\kn{SECMAN: Invalidating
negotiated session rejected by peer}). Restarting Condor on a worker clears its
sessions and re-registers it immediately.

\section{Startup ordering must be enforced, not assumed}

\kn{CONDOR\_HOST} is an mDNS name. At boot, Condor and avahi started \textbf{in
the same second} and it worked only by a hair.

If Condor wins that race, its daemons come up unable to reach their own
collector, and each then sits out a full update interval before retrying. The
pool looks alive and does nothing --- the recurring signature described in
Chapter~\ref{ch:failure}.

\begin{trap}
\kn{After=} alone is insufficient. A \kn{Type=simple} unit is considered started
the moment it execs, which is before the name it publishes has been established.
\end{trap}

Wait for the \textbf{condition} instead, via an \kn{ExecStartPre} that blocks
until the name resolves. Ship it as a \file{condor.service.d/} drop-in ---
editing the packaged unit gets silently reverted on upgrade --- and have the
script \textbf{always exit 0}, so a slow name can never hang a boot.

\section{Do not list a daemon whose package is not installed}

\kn{KBDD} in \kn{DAEMON\_LIST} without \kn{condor-kbdd} present makes the master
log \kn{Cannot execute (errno=2)} and retry every 34 minutes, forever.

Workers whose policy references neither \kn{ConsoleIdle} nor \kn{KeyboardIdle}
--- which, after \S\ref{sec:console}, is all of them --- should not run it at all.
